\clearpage
\renewcommand{\sectioncolor}{alert}
\section{The gaps, and what to buy}
\markboth{The gaps}{}

Five things are missing. In the order I would fill them:

\begin{center}\small
\begin{tabular}{@{}c>{\raggedright\arraybackslash}p{3.4cm}>{\raggedright\arraybackslash}p{5.5cm}>{\raggedright\arraybackslash}p{6.2cm}@{}}
\toprule
\rowcolor{alert!12}
\textbf{\#} & \textbf{\color{alert}Gap} & \textbf{\color{alert}What I would get} & \textbf{\color{alert}Why this one}\\
\midrule
\rowcolor{bio!8}
\textbf{1} & \textbf{Statistical mechanics} &
 \textbf{Chandler}, \textit{Introduction to Modern Statistical Mechanics} (1987) --- or Kardar,
 \textit{Statistical Physics of Particles} (2007) &
 \textbf{If you buy one book, buy Chandler.} $\sim$250 pages, written for chemists, reaches the
 partition function by chapter 3. Kardar is the physicist's book and has the transfer-matrix and
 renormalisation-group material.\\
\textbf{2} & Probability, properly & Grimmett \& Stirzaker, \textit{Probability and Random Processes} &
 Fills \bk{Module\_8}, which is an empty directory.\\
\textbf{3} & \textbf{Markov chains \& MCMC} & Levin \& Peres, \textit{Markov Chains and Mixing Times} &
 Fills \bk{Module\_7}. \textbf{This is the book that explains why sampling $Z$ fails} --- mixing time
 is the quantitative form of the barrier problem you will meet in §6.\\
\textbf{4} & Quantum mechanics & \textbf{Hall}, \textit{Quantum Theory for Mathematicians} --- or
 Sakurai for the physics &
 Hall is the right one for you: it assumes the mathematics and explains the physics, not the other
 way round. It picks up exactly where Axler §7 leaves off.\\
\textbf{5} & Functional analysis & Reed \& Simon vol.~1 &
 Only if you pursue the Tier 3 open-quantum-systems problem of Part~IX. Kuttler ch.~20--24 covers
 most of what comes before.\\
\bottomrule
\end{tabular}
\end{center}

\begin{plain}
Five books, and only the first is urgent. Everything else on the route you already own. The reason
the list is short is that the mathematics you have been building for two years --- linear algebra,
analysis, measure --- \emph{is} the hard part of statistical mechanics. The physics is a thin layer
of interpretation on top of it.
\end{plain}
